% !TEX program = xelatex
% =====================================================================
%  EJEMPLO — póster ubcdposter (A0 vertical, tres columnas)
%  Compila con XeLaTeX:  latexmk main.tex      (ver latexmkrc)
%  Cambia a A1/A2 o apaisado en las opciones del paquete: el diseño
%  se escala solo.
% =====================================================================
\documentclass{article}
\usepackage[a0]{ubcdposter}   % a0 | a1 | a2 ; añade ",landscape" para apaisado
\usepackage[spanish,es-noshorthands,es-nolayout]{babel}

\graphicspath{{../logos/}{logos/}{figures/}}

% ---------------- CABECERA --------------------------------------------
\title{Título del póster en una o dos líneas, claro y directo}
\subtitle{Subtítulo con el mensaje principal del trabajo}
\author{Nombre Apellidos\textsuperscript{1}, Otra Persona\textsuperscript{1,2}, Tercera Autora\textsuperscript{2}}
\ubcdaffiliation{\textsuperscript{1}Unidad, Instituto (Centro), Ciudad\quad
                 \textsuperscript{2}Otro centro, Ciudad}
\ubcdkicker{UNIDAD · INSTITUTO · CENTRO}
\ubcdevent{CONGRESO · CIUDAD · 2027}
% Logos a la derecha de la cabecera (opcional):
% \ubcdtoplogos{\includegraphics[height=110pt]{logo-a.pdf}\hspace{40pt}\includegraphics[height=110pt]{logo-b.png}}

\begin{document}

\ubcdposterheader

% ============================ COLUMNA 1 ==============================
\begin{ubcdcolumn}{0.31}
  \begin{ubcdblock}{01 · CONTEXTO}{Introducción}
    \ubcdlead{Una entradilla de dos o tres líneas que resume el problema y por qué importa.}
    \vspace{24pt}
    Texto del cuerpo. En un póster conviene escribir poco: frases cortas,
    una idea por párrafo y figuras que se entiendan sin leer el texto.
    \begin{itemize}
      \item Primer punto del contexto.
      \item Segundo punto, con una \textbf{palabra clave} en negrita.
      \item Tercer punto.
    \end{itemize}
  \end{ubcdblock}

  \begin{ubcdblock}{02 · OBJETIVOS}{Objetivos}
    \begin{ubcdrows}
      \ubcdrow{01}{Primer objetivo}{Descripción breve del primer objetivo.}
      \ubcdrow{02}{Segundo objetivo}{Descripción breve del segundo objetivo.}
      \ubcdrow{03}{Tercer objetivo}{Descripción breve del tercer objetivo.}
    \end{ubcdrows}
  \end{ubcdblock}

  \begin{ubcdblock}{03 · MÉTODOS}{Métodos}
    \ubcdflow{3}%
    \ubcdstep{01}{Datos}\ubcdarrow
    \ubcdstephl{02}{Modelo}\ubcdarrow
    \ubcdstep{03}{Validación}\par
    \vspace{28pt}
    Describe el procedimiento en pocas líneas: datos de partida, método y
    cómo se ha evaluado frente a la referencia.
  \end{ubcdblock}
\end{ubcdcolumn}\hfill
% ============================ COLUMNA 2 ==============================
\begin{ubcdcolumn}{0.31}
  \begin{ubcdblock}{04 · RESULTADOS}{Resultado principal}
    \ubcdfigure{placeholder.png}{Figura 1. Pie de figura en monoespaciada: qué se ve y cómo leerlo.}
    \vspace{36pt}
    \begin{minipage}[t]{0.31\linewidth}\ubcdstat[ubcdCrimson]{0,94}{F1 · MÉTODO NUEVO}\end{minipage}\hfill
    \begin{minipage}[t]{0.31\linewidth}\ubcdstat{0,71}{F1 · MÉTODO CLÁSICO}\end{minipage}\hfill
    \begin{minipage}[t]{0.31\linewidth}\ubcdstat{×3}{MÁS RÁPIDO}\end{minipage}
  \end{ubcdblock}

  \begin{ubcdblock}[ubcdnuevo]{05 · RESULTADOS}{Comparación con la referencia}
    {\ubcdtablefont\renewcommand{\arraystretch}{1.5}\setlength{\tabcolsep}{14pt}%
    \begin{tabularx}{\linewidth}{@{\hspace{14pt}}X r r}
      \ubcdthead \ubcdth{Conjunto de datos} & \ubcdth{Clásico} & \ubcdth{Nuevo}\\
      \ubcdtalt Conjunto A & 0,71 & \textbf{0,94}\\
      Conjunto B & 0,68 & \textbf{0,90}\\
      \ubcdtalt Conjunto C & 0,80 & \textbf{0,85}\\
    \end{tabularx}}
    \ubcdcaption{Tabla 1. Bloque con la opción ubcdnuevo (borde discontinuo) para destacar.}
  \end{ubcdblock}

  \begin{ubcdblock}{06 · RESULTADOS}{Resultados secundarios}
    \begin{minipage}[t]{0.48\linewidth}
      \ubcdfigure[0.8\linewidth]{placeholder.png}{Figura 2a. Primer panel.}
    \end{minipage}\hfill
    \begin{minipage}[t]{0.48\linewidth}
      \ubcdfigure[0.8\linewidth]{placeholder.png}{Figura 2b. Segundo panel.}
    \end{minipage}\par
    \vspace{24pt}
    Una o dos frases que expliquen qué añaden estas figuras al resultado principal.
  \end{ubcdblock}
\end{ubcdcolumn}\hfill
% ============================ COLUMNA 3 ==============================
\begin{ubcdcolumn}{0.31}
  \begin{ubcdblock}[ubcdpaper]{07 · HERRAMIENTAS}{Software desarrollado}
    \begin{ubcdrows}
      \ubcdtoolrow{Herramienta A}{Qué hace y para quién.}{\ubcdpill[ok]{PÚBLICO}}
      \ubcdtoolrow{Herramienta B}{Qué hace y para quién.}{\ubcdpill[warn]{EN DESARROLLO}}
      \ubcdtoolrow{Herramienta C}{Qué hace y para quién.}{\ubcdpill{BORRADOR}}
    \end{ubcdrows}
  \end{ubcdblock}

  \ubcdband{CONCLUSIÓN}{El mensaje que alguien debería llevarse tras mirar el póster diez segundos.}

  \begin{ubcdblock}{08 · CONCLUSIONES}{Conclusiones}
    \begin{itemize}
      \item Primera conclusión, apoyada en la Figura 1.
      \item Segunda conclusión, apoyada en la Tabla 1.
      \item Siguientes pasos.
    \end{itemize}
  \end{ubcdblock}

  \begin{ubcdblock}{09 · DISPONIBILIDAD}{Código y datos}
    \begin{minipage}[c]{0.62\linewidth}
      Código y datos en repositorios públicos. Escanea el código QR para ver la documentación.\par
      \vspace{12pt}{\ttfamily\small github.com/org/repo\par}
    \end{minipage}\hfill
    \begin{minipage}[c]{0.32\linewidth}\raggedleft
      \ubcdqr[150pt]{https://example.org}
    \end{minipage}
  \end{ubcdblock}

  \begin{ubcdblock}[ubcdplain]{}{Referencias}
    \begin{ubcdrefs}
      \item Autor A, Autor B. Título del artículo. \emph{Revista} 12, 345–356 (2025).
      \item Autor C \emph{et al.} Título del artículo. \emph{Revista} 7, 89 (2026).
    \end{ubcdrefs}
  \end{ubcdblock}
\end{ubcdcolumn}

% ============================== PIE ==================================
\ubcdposterfooter{%
  {\bfseries\color{white}Contacto}\quad nombre.apellidos@centro.es · sitio.web\par
  \vspace{10pt}%
  Agradecimientos y financiación: proyecto, convocatoria o entidad financiadora.}%
 {\ubcdqr{https://example.org}}

% Franja blanca opcional para logos de financiación (según su manual de uso):
% \ubcdposterfunders{\includegraphics[height=120pt]{institucionales.pdf}}

\end{document}